\clearpage
\thispagestyle{frontmatterstyle}
\phantomsection
\label{page:abstract-cn}
\addcontentsline{toc}{chapter}{摘要}

\JOUFrontmatterTitle{摘\hspace{2em}要}

{\zihao{-4}\songti
旋流-静态微泡浮选是一种具有我国自主知识产权的新型柱式分选方法与设备。特有的旋流场结构以及在煤炭分选方面的成功应用，为浮选柱技术在矿物分选领域的进一步拓展奠定了基础。本文围绕浮选旋流分选过程中的流场结构、颗粒受力和分选行为展开分析，构建了研究旋流分选机理的实验与理论框架。

\vspace{0.8\baselineskip}

研究中首先梳理了旋流-静态微泡浮选柱的发展背景与国内外研究现状，在此基础上提出研究目标与技术路线；随后结合流体力学分析、实验测试与典型工况对比，对循环矿浆量、给矿量和柱体背压对分选效果的影响进行了归纳总结；最后从动力学角度讨论旋流场对颗粒分层与界面行为的影响，给出研究结论与后续优化方向。

\vspace{0.8\baselineskip}

本文可为浮选旋流设备的结构优化、参数调节和工程应用提供研究依据，也为后续相关矿物分选装备的设计提供参考。
}

\JOUCNKeywords{浮选；旋流；分选机理；浮选动力学；矿物分选}
